\section{Geometric Unification of Nuclear Forces via Vacuum Saturation}
\label{sec:nukleer_kuvvetler}

In conventional physics, gravitation, electromagnetism, the strong nuclear force, and the weak nuclear force are regarded as four distinct and irreconcilable fundamental interactions. Within the Unified Contraction Hypothesis, however, all these forces emerge as natural manifestations of \textbf{vacuum flow geometry across distinct physical length and density scales}.

\subsection{The Strong Nuclear Force and Quark Confinement (Vacuum Saturation Limit)}
At macroscopic distances, gravitational vacuum inflow dilutes spherically: $v(r) = -\sqrt{2GM/r}$ and $g(r) \propto 1/r^2$. In this regime, the vacuum texture remains unconstrained, and the flux is far below saturation.

However, descending to femtometer scales ($r \sim 10^{-15}\text{ m}$), the immense localized contraction density of quarks draws an inflow that hits the \textbf{maximum carrying capacity of the vacuum (the Planckian saturation threshold)}. The vacuum flux can no longer spread spherically ($1/r^2$ geometric thinning breaks down); the flow between two quarks is channeled into a 1-dimensional cylindrical vortex tube (\textit{flux tube} / vacuum string).

\begin{enumerate}
    \item \textbf{Quark Confinement:}
          Because the cross-section of the vacuum flux tube saturates, the energy per unit length of the tube is constant: $\sigma \approx 1\text{ GeV/fm} \approx 1.6 \times 10^5\text{ N}$ (string tension). The effective potential between two quarks is:
          \begin{equation}
              V_{\mathrm{strong}}(r) = -\frac{4}{3}\frac{\alpha_s \hbar c}{r} + \sigma r.
              \label{eq:guclu_potansiyel}
          \end{equation}
          When attempting to isolate quarks, as the separation $r$ increases, the potential energy $\sigma r$ grows linearly toward infinity. The tension escalates until the vacuum fabric undergoes pair production, yielding a new quark-antiquark pair ($q\bar{q}$); consequently, an isolated free quark can never be observed in the universe!

    \item \textbf{Asymptotic Freedom:}
          When quarks approach one another ($r \to 0$), their contraction basins overlap. The intervening vacuum saturation tube relaxes, and the linear confinement term vanishes ($\sigma r \to 0$). Quarks behave as quasi-free particles within the hadronic boundary.

    \item \textbf{Topological Color Charge and $120^\circ$ Chiral Phase Balance ($SU(3)$ Solution):}
          The three abstract color charges of the Standard Model ($SU(3)$ singlet / ``colorless hadron'' condition) represent, in Scale-Time Calculus, the \textbf{$120^\circ$ ($2\pi/3$) symmetric chiral phase equilibrium} required for three quark vortices ($\mathbb{S}_k = \mathbf{X}_k + \kappa \lambda e^{i \phi_k}, \kappa^2 = 0$) to form a stable closed basin without leaking net rotational vorticity into external space:
          \begin{equation}
              \sum_{k=1}^3 e^{i \phi_k} = 1 + e^{i 2\pi/3} + e^{i 4\pi/3} = 1 + \left( -\frac{1}{2} + i\frac{\sqrt{3}}{2} \right) + \left( -\frac{1}{2} - i\frac{\sqrt{3}}{2} \right) \equiv 0.
              \label{eq:renk_faz_toplami_main}
          \end{equation}
          Analogous to a balanced three-phase alternating current circuit where the net neutral return current vanishes identically, three vortices rotating with relative $120^\circ$ phase offsets cancel each other's tangential vortex drag at the hadron boundary. No net curl leaks into asymptotic space; viewed from outside, the hadron appears completely neutral and ``colorless''. In mesons ($q\bar{q}$), this balance is achieved via opposite chirality and nilpotent superposition ($\kappa^2 = 0$).
\end{enumerate}

\subsection{The Weak Nuclear Force: Vortex Leakage and Chiral Phase Diffraction}
The weak nuclear interaction (e.g., neutron beta decay into a proton, electron, and antineutrino: $n \to p + e^- + \bar{\nu}_e$) represents \textbf{chiral phase leakage} from a contracting nucleon vortex.

During contraction, the internal phase of the quark vortex experiences a minute chirality shear under the extreme saturation pressure of the static vacuum. This leakage is quantized as a localized phase ripple on the vacuum, manifested as massive $W^\pm$ or $Z^0$ vector bosons. The massive $\approx 80\text{ GeV}$ scale of the $W$ boson corresponds to the minimum phase tension required to rupture the saturated nuclear vacuum wall. The extremely short range of the weak interaction ($r_W \approx \hbar / (M_W c) \sim 10^{-18}\text{ m}$) is a geometric suppression enforced by this rigid saturation barrier (Fermi constant $G_F \sim 1/M_W^2$).

\begin{enumerate}
    \item \textbf{Free Neutron Decay ($N=1$):} In an isolated free neutron, lacking an external shielding collective, the internal chiral twist under the Scale-Covariant Derivative $\mathbb{D}_t = \frac{\partial}{\partial t} + \kappa H_0$ overcomes the saturation barrier via WKB tunneling ($\tau_n \approx 879\text{ s} \approx 14.6\text{ minutes}$).
    \item \textbf{Intranuclear Stability (Collective Phase Locking):} When a neutron is bound inside an atomic nucleus, the flux tubes of neighboring nucleons interweave into a shared protective shield ($\mathbb{S}_{\mathrm{nucleus}} = \sum \mathbb{S}_i$). Because the nuclear binding energy per nucleon ($\Delta B \sim 8\text{ MeV/nucleon}$) exceeds the beta decay $Q$-value ($Q \approx 0.782\text{ MeV}$) by an order of magnitude, the nilpotent projection operator $\hat{\mathcal{P}}_\kappa$ spectrally locks the leakage channel. A neutron, which decays in 15 minutes when isolated, achieves absolute stability inside stable nuclei over timescales exceeding the age of the universe ($> 10^{30}\text{ yr}$).
\end{enumerate}

\begin{quote}
    \small
    \textbf{Conceptual and Physical Summary (Effective Geometric Regimes of Forces):} \\
    1. \textbf{Gravitation ($\gamma = 1$ Regime):} The free vacuum inflow potential under spherical $3\mathrm{D}$ dilution at asymptotic macroscopic distances ($v^2 \propto 1/r$, $\Phi \propto 1/r$). \\
    2. \textbf{Strong Interaction ($\gamma = -1$ Regime):} The $1\mathrm{D}$ cylindrical flux tube tension formed when the vacuum carrying capacity reaches its Planckian saturation threshold at femtometer scales ($V(r) \propto \sigma r$). Quark confinement and asymptotic freedom are topological consequences of this saturation geometry. \\
    3. \textbf{Weak Interaction (Chiral Tunneling Regime):} Chiral phase shear and WKB phase leakage across the saturated nuclear vacuum boundary; the effective masses of the $W^\pm$ and $Z^0$ bosons represent the minimum phase tension required to breach this saturation barrier.
\end{quote}

\subsection{Neutrino Oscillations and Chiral Phase Resonance of the PMNS Matrix}
\label{sec:notrino_pmns_main}

In standard particle physics, the quark-sector Cabibbo--Kobayashi--Maskawa (CKM) mixing angles are notably small ($\theta_{12}^{\mathrm{CKM}} = \theta_C \approx 13^\circ$, $\theta_{23}^{\mathrm{CKM}} \approx 2.4^\circ$, $\theta_{13}^{\mathrm{CKM}} \approx 0.2^\circ$), whereas the lepton-sector Pontecorvo--Maki--Nakagawa--Sakata (PMNS) neutrino mixing angles are large and near-maximal ($\theta_{12} \approx 33.4^\circ$, $\theta_{23} \approx 45^\circ$, $\theta_{13} \approx 8.6^\circ$). This stark contrast constitutes an unresolved symmetry puzzle.

Under the Unified Contraction Hypothesis, this hierarchy originates from the \textbf{topological dichotomy between confinement geometry and unconfined phase propagation}:
\begin{enumerate}
    \item \textbf{Quark Sector (CKM / Small Mixing):} Quarks are strictly confined within a 1-dimensional saturated cylindrical flux tube ($\gamma = -1$) on femtometer scales. This rigid saturation geometry restricts the phase shear of quark vortices to a narrow angular band; the Cabibbo angle manifests as a small tangential shear at the nucleon contraction boundary ($\theta_C \approx 13^\circ$).
    \item \textbf{Neutrino Sector (PMNS / Large and Maximal Mixing):} Neutrinos are unconfined, massless or ultra-light chiral phase vortices propagating freely through the static vacuum fabric. Free from hadronic boundary constraints, the three flavor modes ($\nu_e, \nu_\mu, \nu_\tau$) engage in coherent phase resonance with the global chiral background tension of the vacuum over macroscopic distances. The unitary transformation between flavor states and vacuum mass eigenstates is given by:
          \begin{equation}
              |\nu_\alpha\rangle = \sum_{i=1}^3 U_{\alpha i}^{\mathrm{PMNS}} |\nu_i\rangle,
              \label{eq:pmns_donusum}
          \end{equation}
          and the two-flavor transition probability is:
          \begin{equation}
              P(\nu_\alpha \to \nu_\beta) = \sin^2(2\theta) \sin^2\left( \frac{\Delta m^2 c^3 L}{4 \hbar E} \right),
              \label{eq:notrino_salinim}
          \end{equation}
          where under vacuum chiral tension, the resonance naturally reaches the maximal limit for atmospheric neutrinos ($\theta_{23} \approx 45^\circ$) and large mixing for solar neutrinos ($\theta_{12} \approx 33.4^\circ$).
\end{enumerate}
Thus, the contrast between small quark mixing and near-maximal neutrino mixing is resolved as the direct geometric difference between 1D hadronic confinement and free 3D vacuum phase propagation.

\subsection{Force Hierarchy and Density-Driven Dimensional Crossover: Atom Number, Planck Saturation, and Geometric Dimension (\texorpdfstring{$\gamma$}{gamma})}
\label{sec:boyutsal_gecis_main}

In standard physics, the enormous disparity of $\sim 10^{36}$--$10^{40}$ orders of magnitude between gravitation and microscopic interactions is termed the ``Hierarchy Problem''. In the Unified Contraction model, this hierarchy represents a \textbf{density-driven dimensional phase crossover} governed by the ratio of local matter-energy density to the Planck saturation threshold ($c^4/G$):
\begin{equation}
    \left( \frac{v(r)}{c} \right)^2 = 2 \, \mathcal{C}_0 \left( \frac{r}{\ell_0} \right)^{-\gamma}.
    \label{eq:master_formula_main}
\end{equation}
Here, $\mathcal{C}_0$ is the dimensionless fundamental coupling constant, and $\ell_0$ is the characteristic transition length scale (e.g., the Schwarzschild radius $r_s$ in gravitation, or the femtometer scale $\ell_{\mathrm{QCD}}$ in the hadronic sector); this formulation guarantees absolute dimensional consistency across all regimes even as the dimensional index $\gamma$ varies continuously.

The fundamental question arises: \textit{``Is force strength determined by the number of atoms, or by the geometric dimension ($\gamma$)?''}

These two concepts form a causal chain:
\begin{enumerate}
    \item \textbf{The Solar Example (Large $N$, Low Density $\implies \gamma = 1$):} The Sun contains $N \sim 10^{57}$ atoms. Although this immense particle count produces a huge aggregate vacuum inflow ($M = \sum m_i$), these atoms are distributed across an astronomical volume of hundreds of thousands of kilometers. The vacuum flux density per unit area remains trillions of times below the Planckian saturation limit. Because the flow dilutes freely across a spherical $4\pi r^2$ shell, the geometry is 3-dimensional ($\gamma = 1$). Gravitation appears exceptionally weak because its flux dilutes over 3D space.
    \item \textbf{The Proton Example (Small $N=3$, Ultra-High Density $\implies \gamma = -1$):} A proton contains only 3 quarks. In terms of particle count, it is $10^{57}$ times smaller than the Sun; yet these 3 quarks are confined within an ultra-dense volume of $r \sim 10^{-15}\text{ m}$. At this scale, the contraction flux density exceeds the maximum carrying capacity of the vacuum fabric ($\Phi_{\max} = c^4/G$). The vacuum cannot dilute spherically ($1/r^2$ dilution fails); following the path of least action, the dynamical system collapses into a \textbf{1-dimensional cylindrical vortex tube ($\gamma = -1$)}. Because the tube's cross-section is fixed, tension does not diminish with distance ($\sigma \approx 16\text{ tons}$), giving rise to quark confinement.
    \item \textbf{The Weak Force (Saturation Barrier Tunneling $\implies \gamma \to \text{exponential}$):} The chiral phase leakage ($W^\pm, Z^0$) tunneling out of this saturated nuclear vortex is exponentially attenuated by the thickness of the saturation wall ($\mathcal{C}_W \propto e^{-r/r_W}$).
\end{enumerate}
\textbf{When local energy density exceeds Planckian saturation, the flow undergoes a dimensional crossover from free 3D sphericity ($\gamma = 1$) to a 1D string-like flux tube ($\gamma = -1$).}

\subsubsection{Continuous Saturation Phase Crossover for the Dimensional Index \texorpdfstring{$\gamma$}{gamma}}
\label{sec:surekli_gamma}
In our model, the index $\gamma$ is a continuous function of energy density governing the fluid-dynamical saturation crossover:
\begin{equation}
    \gamma(\rho) = \tanh\left( \ln\left[ \frac{\rho_{\mathrm{sat}}}{\rho(r)} \right] \right),
    \label{eq:gamma_surekli_gecis}
\end{equation}
where $\rho_{\mathrm{sat}} \approx 10^{18}\text{ kg/m}^3$ is the hadronic nuclear saturation density.
\begin{itemize}
    \item \textbf{Low-Density Regime ($\rho \ll \rho_{\mathrm{sat}}$):} In macroscopic gravitation and atomic electron clouds, $\frac{\rho_{\mathrm{sat}}}{\rho} \gg 1$, yielding $\ln(\rho_{\mathrm{sat}}/\rho) \to +\infty$ and $\tanh \to +1$. Thus, $\gamma \equiv +1$ (free 3D spherical flux) is recovered precisely.
    \item \textbf{Ultra-Dense Regime ($\rho \gg \rho_{\mathrm{sat}}$):} At the quark scale ($r \le 10^{-15}\text{ m}$), local density exceeds Planck saturation, leading to $\frac{\rho_{\mathrm{sat}}}{\rho} \ll 1$, $\ln(\rho_{\mathrm{sat}}/\rho) \to -\infty$, and $\tanh \to -1$. Consequently, $\gamma \equiv -1$ (1D confined flux tube / quark confinement) emerges smoothly.
\end{itemize}

\subsubsection{Resolution of the Hierarchy Problem: Algebraic Cancellation of $G$ and Cornell Tension}
\label{sec:hiyerarsi_monograph_en}
Addressing why the string tension is $\sigma \approx 1.6 \times 10^5\text{ N}$ ($\approx 1\text{ GeV/fm}$) while the vacuum saturation ceiling is $F_{\mathrm{Planck}} = c^4/G \approx 1.21 \times 10^{44}\text{ N}$, the bridge is provided by the dimensionless quantum gravitational coupling $\alpha_G \equiv G m_q^2 / \hbar c$:
\begin{equation}
    \sigma = \alpha_G \cdot F_{\mathrm{Planck}} = \left( \frac{G m_q^2}{\hbar c} \right) \left( \frac{c^4}{G} \right) = \frac{m_q^2 c^3}{\hbar}.
    \label{eq:sigma_monograph_ispat_en}
\end{equation}
Remarkably, Newton's gravitational constant $G$ cancels identically. For a constituent quark mass $m_q \approx 440\text{ MeV}$, this yields $\sigma \approx 1.57 \times 10^5\text{ N} \approx 0.98\text{ GeV/fm}$, resolving the 39-order hierarchy without arbitrary parameter tuning.

\subsection{Epistemological Foundation: Reduction of Four Fundamental Postulates into a Single Parameter}
\label{sec:postulatlarin_eritilmesi_sentez}

The historical failure of unified field theories over the past century stems from attempting to unify the four interactions while retaining \textbf{time ($t$), position ($\mathbf{x}$), mass ($m$), and the speed of light ($c$)} as disjoint, mutually independent postulates.

To unify the four fundamental interactions into a single dynamical framework, these four quantities must themselves be derived as geometric and kinematic projections of a single underlying mechanism: the scale contraction dynamics of matter within the vacuum ($\lambda(t)$ and $\kappa$). This ontological reduction is summarized in Table~\ref{tab:ontolojik_indirgeme_sentez}.

\begin{table}[h!]
    \centering
    \small
    \begin{tabular}{@{}lll@{}}
        \toprule
        \textbf{Classical Postulate} & \textbf{Equivalent in Scale-Time Calculus}                                                                         & \textbf{Physical Mechanism}                              \\ \midrule
        \textbf{Speed of Light ($c$)} & $v_c = 1/\sqrt{\varepsilon_0\mu_0} = c$                                                                            & Elastic vacuum phase propagation speed                   \\
        \textbf{Time ($t$)}           & $dt \equiv -\frac{1}{H_0} \frac{d\lambda}{\lambda} = -\frac{1}{H_0} d(\ln\lambda)$                                 & Quantum phase log recorded by contracting atoms          \\
        \textbf{Position ($R$)}       & $\mathbf{R}(t) \equiv \frac{\mathbf{X}_0}{\lambda(t)} = \mathbf{X}_0 e^{H_0 t}$                                    & Ratio of static space to the contracting local ruler     \\
        \textbf{Mass ($M$)}           & $M \equiv \frac{1}{4\pi G} \oint_{\partial \Omega} |(\mathbf{v}_{\mathrm{inflow}}\cdot\nabla)\mathbf{v}_{\mathrm{inflow}}| \, dA$ & Net convective acceleration flux of vacuum inflow        \\ \bottomrule
    \end{tabular}
    \caption{Reduction of the four classical postulates into the single scale-contraction dynamic ($\lambda(t)$).}
    \label{tab:ontolojik_indirgeme_sentez}
\end{table}

With these four concepts rooted in a single mechanism, the flow dynamics of Gravitation, Electromagnetism, the Strong Force, and the Weak Force can be classified and modeled as geometric manifestations of the same vacuum inflow ($\mathbf{v}_{\mathrm{inflow}}$) across varying density and saturation regimes.

\subsection{Geometric Synthesis: Unified Inflow Dynamics of the Four Interactions}
\label{sec:ana_dinamik_denklem_main}

Combining the kinematic acceleration relation of Scale-Time Calculus in observer coordinates with the Unified Master Inflow Formula (\ref{eq:master_formula_main}) yields the \textbf{Master Dynamic Equation} governing microscopic and macroscopic systems within a single framework:
\begin{equation}
    \boxed{ m \frac{d^2\mathbf{R}}{dt^2} = - m c^2 \nabla \left[ \mathcal{C}_0 \left(\frac{R}{\ell_0}\right)^{-\gamma} \right] + 2 m H_0 \mathbf{v}_{\mathrm{local}} + m H_0^2 \mathbf{R} }
    \label{eq:master_dynamic_equation}
\end{equation}

This single differential equation accounts for all operational regimes of the cosmos:
\begin{enumerate}
    \item \textbf{First Term (Vacuum Inflow Potential of the Four Forces):} The term $- m c^2 \nabla [\mathcal{C}_0 (R/\ell_0)^{-\gamma}]$ adapts according to the local matter-energy density:
          \begin{itemize}
              \item $\gamma = 1 \implies$ Gravitation ($\mathcal{C}_0 = GM/(c^2 \ell_0)$) and Electromagnetism ($\mathcal{C}_0 = \alpha \hbar / (m_e c \ell_0)$).
              \item $\gamma = -1 \implies$ Strong Nuclear Force / Quark Confinement ($\mathcal{C}_0 = \sigma \ell_0 / (2m_q c^2)$).
              \item $\gamma \to \text{exponential} \implies$ Weak Nuclear Force / Beta Decay ($\mathcal{C}_0 \propto e^{-R/r_W}$).
          \end{itemize}
    \item \textbf{Second Term (Scale Phase Damping):} $2 m H_0 \mathbf{v}_{\mathrm{local}}$ is the cosmic Coriolis drag experienced by motion within the contracting reference frame.
    \item \textbf{Third Term (Cosmic Expansion Apparition / Dark Energy Equivalent):} $m H_0^2 \mathbf{R}$ represents the geometric acceleration that standard cosmology attributes to dark energy ($\Lambda \approx 3 H_0^2$), arising naturally from the differential calculus of shrinking coordinates.
\end{enumerate}

Thus, the theory achieves an exact analytical synthesis at two complementary levels:
\begin{itemize}
    \item \textbf{Kinematic Inflow Field:} $\left(\frac{v(r)}{c}\right)^2 = 2 \, \mathcal{C}_0 \left(\frac{r}{\ell_0}\right)^{-\gamma}$ (Static Vacuum Inflow).
    \item \textbf{Dynamical Trajectory Equation:} $m \frac{d^2\mathbf{R}}{dt^2} = - m c^2 \nabla \left[ \mathcal{C}_0 \left(\frac{R}{\ell_0}\right)^{-\gamma} \right] + 2 m H_0 \mathbf{v} + m H_0^2 \mathbf{R}$ (Scale-Covariant Motion).
\end{itemize}
